Agent memory has emerged as a first-class primitive in the pursuit of intelligent, long-horizon AI agents. The original survey by \citet{DBLP:journals/corr/abs-2512-13564} established a unified $3\times 3 \times 3$ taxonomy---Forms $\times$ Functions $\times$ Dynamics---organizing roughly 200 papers up to January 2026. Since then, the field has accelerated dramatically: our living catalog now spans 952 papers (as of July 2026), including 498 works published since February 2026---an eightfold increase over the original survey's post-cutoff additions. This explosive growth reveals cross-cutting themes the original taxonomy was not designed to capture. In this paper we make three contributions. First, we describe a fully data-driven methodology for curating and maintaining a living paper survey, with a single-source-of-truth YAML file, automated validation pipelines, bulk metadata enrichment, and an open-source infrastructure at \url{github.com/tobias-weiss-ai-xr/agent-memory-research}. Second, we extend the original taxonomy with three novel orthogonal dimensions---Temporal Dynamics, Modality, and Biological Inspiration---that capture the emerging themes of forgetting curves, multimodal memory, and neurocognitively inspired architectures, and we document two further theme clusters---security and adversarial robustness, and efficiency-driven compression---that have grown into major research strands in their own right. Third, we conduct a gap analysis that identifies persistently underpopulated taxonomy cells, missing evaluation standards, and a concrete research roadmap spanning short, medium, and long-term milestones. All data, code, and a browsable web interface accompany this paper.
